\documentclass[11pt]{letter}
\usepackage[margin=2.5cm]{geometry}

\begin{document}

\begin{letter}{The Editors \\ \textit{Communications in Mathematical Physics}}

\opening{Dear Editors,}

We submit for your consideration our manuscript entitled
\textbf{``Gap ratio statistics of Riemann zeros: measurement, mechanism,
and the Berry-Keating correction''}
for publication in \textit{Communications in Mathematical Physics}.

This work presents the first precision measurement of the rate at which
the gap ratio statistic of Riemann zeta zeros converges to the GUE
prediction of random matrix theory. Using 21 independent data points
spanning $\log T = 9.7$ to $24.1$, we find
$\langle r\rangle(T) = 0.59891(13) + 1.245(40)/\log^2 T$,
with the asymptotic value $6.1\sigma$ below the GUE limit.

We identify the physical mechanism producing the correction: the spacing
distribution of Riemann zeros is narrower than GUE (contributing $+128\%$
to the coefficient $c$) while anti-correlation is stronger ($-29\%$),
reproducing $99.5\%$ of the measured value. Independent confirmation
comes from the number variance $\Sigma^2$ and spectral rigidity
$\Delta_3$, which exhibit the Bogomolny--Keating saturation.

Two companion papers develop related results: the ab initio derivation
of $c$ from the Conrey--Snaith triple correlation formula (Paper~B),
and the unconditional proof that Rudnick--Sarnak correlations and
linear programming imply the Riemann Hypothesis (Paper~C).

We believe this work is appropriate for \textit{CMP} because it
establishes a precise quantitative link between prime number distribution
and random matrix theory, connecting empirical measurement with
analytical mechanism via Painlev\'e~V theory.

We confirm that this manuscript has not been published previously and is
not under consideration elsewhere.

\closing{Sincerely,}

David Alarc\'on\\
Universidad Pablo de Olavide, Sevilla, Spain\\
\texttt{dalarub@upo.es}

\end{letter}
\end{document}
